\section{Setup and experimental details}

\subsection{Pew research surveys}
\label{app:pew_collection}

Our dataset is derived from the annual Pew American Trends Panel (ATP) survey.
Below, we provide a brief summary of how the data collection process is conducted, and refer the reader to \url{pewresearch.org/our-methods/u-s-surveys/the-american-trends-panel/} and \url{pewresearch.org/our-methods/u-s-surveys/writing-survey-questions/} for more details.


\paragraph{Panelists.}

For ATP surveys, Pew relies on a group of about 10,000 participants within the US recruited over multiple years, many of whom take the survey repeatedly.
Each year, a subset of panelists are invited to take the ATP to reduce the burden on individual respondents.
Panelists are offered a paid incentive to participate in the survey.

Panelists are recruited by sending participation requests to a randomly-chosen address-based sample of households from USPS's Delivery Sequence File with concerted efforts to ensure representativeness of the sample. 
They also solicit input from households without internet access---either via phone or by providing them with tablets to take the survey.

\paragraph{Questionairre design.}
As stated on the Pew research website: "Perhaps the most important part of the survey process is the creation of questions that accurately measure the opinions, experiences and behaviors of the public...Designing the questionnaire is complicated because surveys can ask about topics in varying degrees of detail, questions can be asked in different ways, and questions asked earlier in a survey may influence how people respond to later questions."

Pew research selects pertinent topics for their surveys by monitoring the state of the nation and the world, and identifying issues that would be relevant to the public, media and policymakers.
They then go through an iterative process to build questions, often piloting them in focus groups, pre-interviews and cognitive testing.
The question wording is highly optimized to be clear, easy-to-understand, and not bias participants towards a particular answer. 

In order to identify valid choices for questions, Pew researchers often initially pilot open-ended surveys, and then use them to determine valid answer choices.

\paragraph{Data quality.}
Every survey, once designed is first tested out on a set of 60 ``fast'' panelists to flag any design errors.
Pew researchers also conduct data quality checks to identify issues with respondent satisfaction or the collected answers.
The ATP data is also accompanied with sample weights per individual to account for sampling bias and non-response over various stages of data collection.

Researchers have observed that human participants are sensitive to question and option ordering.
However, for questions with ordinal options ("Strongly agree"..."Strong disagree"), the option ordering is not randomized since they view it as conveying important information.

\subsection{Adapting ATP to \bname}
\label{app:survey_curation}
\label{app:taxonomize}

We derive our questions and human reference distributions based on \nsurveys ATP surveys over multiple years (2017-2021)---see Appendix Table~\ref{tab:survey_stats} for details.
The prefix in each survey name points to the wave in which it was collected.
We chose these surveys as they span a broad range of topics that might be pertinent for human-centric LM applications.
In Appendix Table~\ref{tab:survey_md}, we depict the demographic traits that we consider in our sub-group level analysis. 


\paragraph{Post-processing.}
As such, we directly extract multiple-choice questions from Pew ATP surveys and try to apply as little post-processing as possible. Some cases where we must filter or modify the questions are:
\begin{enumerate}
\item \emph{Cross-references:} Some questions make explicit references to context provided in a previous question. However, since we are presenting questions to LMs individually, we must modify every question to be self-contained. 
\item \emph{Variable-dependent questions:} We omit questions  where the phrasing of the question itself depends on a previous answer: ``In your answer to the previous question, you said \$ANSWER. Is this because....''. 
\item \emph{Formatting:} We fix any formatting issues that in questions to make them suitable for LMs (e.g., weird tokens or all capital words).
\item \emph{Lists:} Often, Pew surveys have lists where the same question is asked of many different variables. For instance, ``How much does each of the following affect your happiness in life? [A lot/.../Not at all]'' followed by a series of \$Xs such as ``money'', ``exercise''... In these cases, we restate the question to be self-contained, i.e., ``How much does \$X affect your happiness in life? [A lot/.../Not at all]'' in this case.
\end{enumerate}

As stated above, we try to keep our edits as minimal as possible.
In Appendix Table~\ref{tab:topic}, we describe the categories we manually taxonomize our dataset into for post-hoc topic-level analysis.
Note that questions may fall into multiple categories.

\begin{table*}[!h]
\caption{Summary of Pew surveys used in our analysis: $N_Q$ and $N_R$ denote the number of questions and human respondents respectively. (Continued on next page.)}
\label{tab:survey_stats}
\vskip 0.15in
\renewcommand{\arraystretch}{1.25}
\begin{center}
\begin{small}

    \begin{tabular}{llp{1.25cm}ccp{5cm}}
        \toprule
            Name &      Field dates &                Topic &  \# Questions &  \# Responses &                                                                                                                                                                                                                                                                          Sample question \\
        \midrule
        ATP\_W26  &   April 4-18, 2017      &              Guns &           78 &         4168 &                                                                                                In general, as far as you know, how many of the guns in your home would you say are kept loaded? [All are kept loaded/Some are kept loaded and some are not/None are kept loaded/Refused] \\
        ATP\_W27  &   May 1-15, 2017      &              Automation and driverless vehicles &           96 &         4135 &                                                                                                Would you feel better or worse about computer programs making hiring decisions if these computer programs included public data about each candidate - such as the material they post on social media - in making their evaluations [Better/Worse/No difference/Refused] \\
        ATP\_W29  &   Sept 14–28, 2017  &         Views on gender &           77 &         4867 &                                                                                Thinking about how society sees men these days, in general, would you say [Most people look up to men who are manly or masculine/Most people look down on men who are manly or masculine/Neither/Refused] \\
        ATP\_W32  & Feb 26–March 11, 2018 &          Community types and sexual harassment  &           98 &         6251 &                                                                                                                                                                                                    How important is it to you, personally, to live in a community that is a good place to raise children [Very important/Somewhat important/Not too important/Not at all important/Refused] \\
        \bottomrule
        \end{tabular}
    \end{small}
\end{center}    
\vskip -0.1in
\end{table*}

\begin{table*}[t]
    \ContinuedFloat
    \caption{Summary of Pew surveys used in our analysis: $N_Q$ and $N_R$ denote the number of questions and human respondents respectively.}
    \vskip 0.15in
    \renewcommand{\arraystretch}{1.25}
    \begin{center}
    \begin{small}
    
        \begin{tabular}{lcp{1.25cm}ccp{5cm}}
            \toprule
                Name &      Time period &                Topic &  \# Questions &  \# Responses &                                                                                                                                                                                                                                                                          Sample question \\
            \midrule
            ATP\_W34  & April 26–May 6, 2018 &   Biomedical and food issues &           67 &         2537 &                                    In your opinion, do you think government investments in engineering and technology usually pay off in the long run, or are they not worth it? [Government investments usually pay off in the long run/Government investments aren't worth it/Refused] \\
            ATP\_W36  & June 19–July 2, 2018 &      Gender and leadership &          139 &         4587 &                                                                                                                 In general, do you think men or women in top executive business positions are better at working out compromises? [Men are better/Women are better/No difference/Refused] \\
            ATP\_W41  & Dec 10–23, 2018  &           America in 2050 &           90 &         2524 &                                            In the future, what kind of an impact do you think the news media will have in solving the biggest problems facing the country? [A very positive impact/A somewhat positive impact/A somewhat negative impact/A very negative impact/Refused] \\
            ATP\_W42  & Jan 7–21, 2019  &          Trust in science &          129 &         4464 &                                                                                                    When you hear or read news stories about research misconduct by nutrition research scientists, do you think of these cases as [Isolated incidents/Signs of a broader problem/Refused] \\
            ATP\_W43  &  Jan 22–Feb 5, 2019   &                     Race &          114 &         6637 & For each, please indicate if you, personally, think it is acceptable. A white person using makeup to darken their skin so they appear to be a different race as part of a Halloween costume [Always acceptable/Sometimes acceptable/Rarely acceptable/Never acceptable/Not sure/Refused] \\
            ATP\_W45  &    Feb 19–March 4, 2019 &           Misinformation &           95 &         6127 &                                                                                                                                                                          How much made-up news and information do you think is created by journalists [A lot/Some/Not much/None/Refused] \\
            ATP\_W49  &  June 3–17, 2019 &    Privacy and surveillance &           98 &         4272 &                                                                                                                                       How much do you feel you understand what companies are doing with the data they collect about you? [A great deal/Some/Very little/Nothing/Refused] \\
        \bottomrule
        \end{tabular}

\end{small}
\end{center}    
\vskip -0.1in
\end{table*}

\begin{table*}[t]
    \ContinuedFloat
    \caption{Summary of Pew surveys used in our analysis: $N_Q$ and $N_R$ denote the number of questions and human respondents respectively.}
    \vskip 0.15in
    \renewcommand{\arraystretch}{1.25}
    \begin{center}
    \begin{small}
    
        \begin{tabular}{lcp{1.25cm}ccp{5cm}}
            \toprule
                Name &      Time period &                Topic &  \# Questions &  \# Responses &                                                                                                                                                                                                                                                                          Sample question \\
            \midrule
        ATP\_W50  & June 25–July 8, 2019 &   Relationships and family &          128 &         9834 &                                                                                                                                               How much, if at all, do you trust your spouse/partner to handle money responsibly [A great deal/A fair amount/Not much/Not at all/Refused] \\
        ATP\_W54  & Sept 16–29, 2019  &       Economic inequality &          116 &         6878 &                                                                                  Do you think the country's current economic conditions are helping or hurting people who are white? [Helping a lot/Helping a little/Hurting a little/Hurting a lot/Neither helping nor hurting/Refused] \\
        ATP\_W82  & Feb 2–7, 2021  &          Global attitudes &          104 &         2596 &               When it comes to whether or not to limit Chinese students studying in the U.S., do you [Strongly support limiting Chinese students/Somewhat support limiting Chinese students/Somewhat oppose limiting Chinese students/Strongly oppose limiting Chinese students/Refused] \\
        ATP\_W92  &  July 8–18, 2021   &         Political views &           77 &        10221 &                                                                                                                      Do you think a decline in the share of Americans belonging to an organized religion is generally good or bad for our society? [Very good for society/Somewhat good for society/Neither good nor bad for society/Somewhat bad for society/Very bad for society/Refused] \\
        \bottomrule
        \end{tabular}

\end{small}
\end{center}    
\vskip -0.1in
\end{table*}
\begin{table*}

   \caption{Summary of demographic traits used in our group-level analysis.}
   \label{tab:survey_md}
   \vskip 0.15in
   \renewcommand{\arraystretch}{1.5}
   \begin{center}
   \begin{small}
   
    \begin{tabular}{lp{5cm}p{8cm}}
        \toprule
        Attribute &                                                                       Interpretation &                                                                                                                                                                                                                                                           options \\
     \midrule
         CREGION &                     Which part of the United States do you currently live in? &                                                                                                                                                                                                                                 [Northeast, Midwest, South, West] \\
          SEX &              What is the sex that you were assigned at birth? &                                                                                                                                                                                                                      [Male, Female] \\
             AGE &                                                              How old are you? &                                                                                                                                                                                                                               [18-29, 30-49, 50-64, 65+] \\
       EDUCATION &     What is the highest level of schooling or degree that you have completed? &                                                                                                                [Less than high school, High school graduate, Some college, no degree, Associate's degree, College graduate/some postgrad, Postgraduate] \\
            RACE &                                                  What is your race or origin? &                                                                                                                                                              [White, Black, Asian, Hispanic, 'Other] \\
         CITIZEN &                                       Are you a citizen of the United States? &                                                                                                                                                                                                                                                [Yes, No] \\
         MARITAL &                                            Which of these best describes you? &                                                                                                                                                                       [Married, Living with a partner, Divorced, Separated, Widowed, Never been married] \\
           RELIG &                                        What is your present religion, if any? &            [Protestant, Roman Catholic, Mormon, Orthodox, Jewish, Muslim, Buddhist, Hindu, Atheist, Agnostic, Other, Nothing in particular] \\
     RELIGATTEND & Aside from weddings and funerals, how often do you attend religious services? &                                                                                                                                                           [More than once a week, Once a week, Once or twice a month, A few times a year, Seldom, Never] \\
        POLPARTY &                                 In politics today, do you consider yourself a &                                                                                                                                                                                                      [Republican, Democrat, Independent, Something else] \\
          INCOME &  Last year, what was your total family income from all sources, before taxes? & [Less than \$30,000, \$30,000-\$50,000, \$50,000 -\$75,000, \$75,000-\$100,000,  \$100,000 or more] \\
     POLIDEOLOGY &                        In general, would you describe your political views as &                                                                                                                                                                                       [Very conservative, Conservative, Moderate, Liberal, Very liberal] \\
     \bottomrule
    \end{tabular}
   \end{small}
\end{center}    
\vskip -0.1in
\end{table*}

\begin{table*}[!h]

    \caption{Topic breakdown of questions in \bname{}; high-level topics are in bold and sub-categories are italicized. Note: a questions can belong to multiple topics. }
    \label{tab:topic}
    \renewcommand{\arraystretch}{1.5}
    \begin{center}
    \begin{small}

        \begin{tabular}{p{4.5cm}rp{10.5cm}}
            \toprule
            Topic &  $N_Q$ & Example \\
            \midrule
            \textbf{community health} &    67 & How important is it to you, personally, to live in a community where most people share your religious views [Very important/Somewhat important/Not too important/Not at all important/Refused]  \\ \hline
            \textbf{corporations, tech, banks and automation} &    107 & \\ 
            \textit{robots} & 43 & Please consider the following scenario - in the future, robots and computers with advanced capabilities may be able to do most of the jobs that are currently done by humans today. How much have you heard, read, or thought about this idea before today? [A lot/A little/Nothing at all/Refused] \\
            \textit{voice assistants} & 7 & When you use digital assistants, how often do they accurately respond to your commands? [Most of the time/Some of the time/Not very often/Refused] \\
            \textit{drones} & 7 & Do you think that private citizens should or should not be allowed to pilot drones in the following areas? Near crime scenes or traffic accidents [Should be allowed/Should not be allowed/It depends/Refused] \\
            \textit{autonomous vehicles} & 17 & How enthusiastic are you, if at all, about the development of driverless vehicles? [Very enthusiastic/Somewhat enthusiastic/Not too enthusiastic/Not at all enthusiastic/Refused] \\
            \textit{other} & 33 & How much power and influence do you think technology companies have on today's economy? [Too much power and influence/Not enough power and influence/About the right amount/Refused] \\ \hline
            \textbf{crime/security} &    89 &  \\
            \textit{crime} & 5 & How much, if at all, do you worry about the following happening to you? Being the victim of a mass shooting [Worry a lot/Worry a little/Do not worry at all/Refused] \\
            \textit{guns} & 73 & Thinking about gun owners who do not have children in their home how important do you think it is for them to: Advise visitors with children that there are guns in the house [Essential/Important but not essential/Not important/Should not be done/Refused] \\  
            \textit{justice system} & 4 & Overall, would you say people who are convicted of crimes in this country serve [Too much time in prison/Too little time in prison/About the right amount of time in prison/Refused] \\
            \textit{military} & 3 & How much confidence, if any, do you have in the military to act in the best interests of the public? [A great deal of confidence/A fair amount of confidence/Not too much confidence/No confidence at all/Refused] \\
            \textit{terrorism} & 5 & Thinking about long-range foreign policy goals, how much priority, if any, do you think taking measures to protect the U.S. from terrorist attacks should be given? [Top priority/Some priority/No priority/Refused] \\  
           \bottomrule
        \end{tabular}
\end{small}
\end{center}    
\vskip -0.1in
\end{table*}
    
\begin{table*}[!h]
    \renewcommand{\arraystretch}{1.5}
    \begin{center}
    \begin{small}
        \begin{tabular}{p{4.5cm}rp{10.5cm}}
            \toprule
            Topic &  $N_Q$ & Example \\
            \midrule
            \textbf{discrimination} &    62 &  \\ 
            \emph{racial} &    36 & Would you say that black people are treated less fairly than white people, white people are treated less fairly than black people, or both are treated about equally in in stores or restaurants situations? [Black people are treated less fairly than white people/White people are treated less fairly than black people/Both are treated about equally/Refused] \\ 
            \emph{sexual harassment} &    21 & When it comes to sexual harassment in the workplace today, how much of a problem, if at all, would you say women claiming they have experienced sexual harassment or assault when it hasn't actually occurred is? [Major problem/Minor problem/Not a problem/Refused]  \\ 
            \emph{other} &    5 & Have you personally experienced the following at work because you have children? Being passed over for a promotion [Yes, have experienced this/No, have not experienced this/Refused] \\  \hline
            \textbf{economy and inequality} &    94 & How much, if at all, do you think not enough regulation of major corporations contributes to economic inequality in this country? [Contributes a great deal/Contributes a fair amount/Contributes not too much/Contributes not at all/Refused] \\ \hline
            \textbf{education} &    27 & Do you think scores on standardized tests, such as the SAT or act should be a major factor, minor factor, or not a factor in college admissions? [Major factor/Minor factor/Not a factor/Refused] \\ \hline
            \textbf{future} &    55 & Thinking again about the year 2050, or 30 years from now, do you think abortion will be [Legal with no restrictions/Legal but with some restrictions/Illegal except in certain cases/Illegal with no exceptions/Refused] \\ \hline
            \textbf{gender \& sexuality} &   165 &  \\
            \textit{gender attitudes} &   155 &  In general, do you think men or women in high political offices are better at standing up for what they believe in, despite political pressure? [Men are better/Women are better/No difference/Refused] \\
            \textit{sexuality} &   10 &  Do you think greater social acceptance of people who are transgender (people who identify as a gender that is different from the sex they were assigned at birth) is generally good or bad for our society? [Very good for society/Somewhat good for society/Neither good nor bad for society/Somewhat bad for society/Very bad for society/Refused]\\ 
            \bottomrule
            \end{tabular}
    \end{small}
 \end{center}    
 \vskip -0.1in
 \end{table*}

 \begin{table*}[!h]
    \ContinuedFloat

    \caption{Topic breakdown of questions in \bname{}; high-level topics are in bold and sub-categories are italicized. Note: a questions can belong to multiple topics. }    \vskip 0.15in
    \renewcommand{\arraystretch}{1.5}
    \begin{center}
    \begin{small}

        \begin{tabular}{lrp{10cm}}
            \toprule
            Topic &  $N_Q$ & Example \\
            \midrule
            \textbf{global attitudes and foreign policy} &    78 & Thinking about long-range foreign policy goals, how much priority, if any, do you think limiting the power and influence of North Korea should be given? [Top priority/Some priority/No priority/Refused]\\
            \textbf{healthcare} &    58 &  \\ 
            \textit{abortion} &    4 & Which statement comes closer to your own views? [There are some situations in which abortion should be allowed/There are no situations at all where abortion should be allowed/Refused] \\ 
            \textit{covid} &    7 & Thinking about restrictions on public activity in the US over the course of the coronavirus outbreak, do you think there should have been [More restrictions/Fewer restrictions/The restrictions were about right/Refused] \\ 
            \textit{other} &    47 & Thinking about medical treatments these days, how much of a problem, if at all, are the following? Healthcare providers are too quick to order tests and procedures that may not be necessary [A big problem/A small problem/Not a problem/Refused] \\ \hline
            \textbf{immigration} &    19 & How much, if at all, do you think the growing number of illegal immigrants working in the U.S. contributes to economic inequality in this country? [Contributes a great deal/Contributes a fair amount/Contributes not too much/Contributes not at all/Refused] \\ \hline
            \textbf{job/career} &    67 & How much, if at all, do you worry about the following happening to you? Losing your job [Worry a lot/Worry a little/Do not worry at all/Refused] \\ \hline
            \textbf{leadership} &    31 & In general, how important, if at all, is it to you for someone in a top executive business position to do be compassionate and empathetic? [Essential/Important, but not essential/Not important/Refused] \\ \hline
            \textbf{news, social media, data, privacy} &   198 & \\
            \textit{data \& privacy} &   85 & Do you think it is possible to go about daily life today without having the government collect data about you? [Yes, it is possible/No, it is not possible/Refused] \\
            \textit{news \& social media} &   113 & How much of a problem is the amount of made-up news and information when it comes to how the public stays informed about the basic facts of current issues and events? [A very big problem/A moderately big problem/A small problem/Not a problem at all/Refused] \\ \hline
            \textbf{personal finance} &    45 & How often, if ever, do you worry about the amount of debt you have? [Every day/Almost every day/Sometimes/Rarely/Never/Refused] \\ \hline
            \textbf{personal health} &    29 & Do you think organic fruits and vegetables are generally [Better for one's health than conventionally grown foods/Worse for one's health than conventionally grown foods/Neither better nor worse for one's health than conventionally grown foods/Refused] \\       \hline
            \textbf{political issues} &   112 &  \\ 
            \textit{Two party system} &   34 &  Since President Trump was elected, do you think it has become more common or less common for people to express racist or racially insensitive views, or is it about as common as it was before? [More common/Less common/About as common/Refused] \\ 
            \textit{government control} &   69 &  Should health insurance [Be provided through a single national health insurance system run by the government/Continue to be provided through a mix of private insurance companies and government programs/Refused] \\
            \textit{fair elections} &   6 &  Still thinking about elections in the country, how confident, if at all, are you that people who are not legally qualified to vote are prevented from casting a ballot [Very confident/Somewhat confident/Not too confident/Not at all confident/Refused] \\       \bottomrule
        \end{tabular}
\end{small}
\end{center}    
\vskip -0.1in
\end{table*}

\begin{table*}[!h]
    \ContinuedFloat

    \caption{Topic breakdown of questions in \bname{}; high-level topics are in bold and sub-categories are italicized. Note: a questions can belong to multiple topics. (Continued on next page)}    \vskip 0.15in
    \renewcommand{\arraystretch}{1.5}
    \begin{center}
    \begin{small}
   
        \begin{tabular}{lrp{10cm}}
            \toprule
            Topic &  $N_Q$ & Example \\
            \midrule
            \textbf{race} &    116 & How much more, if anything, needs to be done to ensure equal rights for all Americans regardless of their racial or ethnic backgrounds? [A lot/A little/Nothing at all/Refused]\\ \hline
            \textbf{relationships and family} &   114 & Looking ahead, would having children make it [Easier to advance in your job or career/Harder to advance in your job or career/Would not make a difference/Refused] \\ \hline
            \textbf{religion} &    12  & Do you think a decline in the share of Americans belonging to an organized religion is generally good or bad for our society? [Very good for society/Somewhat good for society/Neither good nor bad for society/Somewhat bad for society/Very bad for society/Refused]\\ \hline
            \textbf{science} &   160 & Do you think genetic engineering of animals to grow organs or tissues that can be used for humans needing a transplant would be [An appropriate use of technology/Taking technology too far/Refused]\\
            \textit{climate} &   41 & How confident are you, if at all, that the actions taken by the international community will significantly reduce the effects of global climate change? [Very confident/Somewhat confident/Not too confident/Not at all confident/Refused]\\
            \textit{other} &   119 & Do you think genetic engineering of animals to grow organs or tissues that can be used for humans needing a transplant would be [An appropriate use of technology/Taking technology too far/Refused]\\ \hline
            \textbf{self-perception and values} &    40 & How well, if at all, do the following words or phrases describe you? Physically strong [Very well/Somewhat well/Not too well/Not at all well/Refused]\\ \hline
            \textbf{status in life} &    20 & Generally, how would you say things are these days in your life? Would you say that you are [Very happy/Pretty happy/Not too happy/Refused] \\ \hline
            \bottomrule
            \end{tabular}
    \end{small}
 \end{center}    
 \vskip -0.1in
 \end{table*}
 
\begin{table*}

    \caption{Demographic groups used in our steerability analysis.}
    \label{tab:steer_groups}
    \vskip 0.15in
    \renewcommand{\arraystretch}{1.5}
    \begin{center}
    \begin{small}

        \begin{tabular}{ll}
            \toprule
    Attribute &                                            Demographic group \\
            \midrule
            CREGION &                                 Northeast, South \\
            EDUCATION &  College graduate/some postgrad, Less than high school \\
            GENDER &                                     Male, Female \\
            POLIDEOLOGY &                  Liberal, Conservative, Moderate \\
            INCOME &              \$100K+, $<$\$30,000 \\
            POLPARTY &                             Democrat, Republican \\
            RACE &                    Black, White, Asian, Hispanic \\
            RELIG &       Protestant, Jewish, Hindu, Atheist, Muslim \\
            \bottomrule
            \end{tabular}
    

    \end{small}
 \end{center}    
 \vskip -0.1in
 \end{table*}
 

For our steerability analysis in Section~\ref{sec:steerable}, we pick a subset of 500 questions where the subgroups under consideration frequently disagree.

\clearpage



\subsection{Models}
\label{app:models}
For our analysis, we use a series of models from OpenAI and AI21 labs, detailed in \ref{tab:models}.
Since the model training process is not always publicly known, we attempt to report this to the best of our knowledge.
Further documentation can be found \url{beta.openai.com/docs/model-index-for-researchers} and \url{docs.ai21.com/docs/}.
\begin{table*}

    \caption{LLMs we evaluate in our study. In some cases, we attempt to report size/training details of models to the best of our ability as these are often not clearly disclosed.}
    \label{tab:models}
    \vskip 0.15in
    \renewcommand{\arraystretch}{1.5}
    \begin{center}
    \begin{small}

        \begin{tabular}{lllp{8cm}}
            \toprule
    Model name &  Provider & Size & Notes \\
            \midrule
            j1-Grande & AI21 Labs & 17B & Auto-regressive model from \citet{lieber2021jurassic} \\
            j1-Jumbo  & AI21 Labs & 178B & Auto-regressive model from \citet{lieber2021jurassic} \\
            j1-Grande v2 beta & AI21 Labs & 17B &  Instruct tuned version of j1-Grande, trained specifically to handle zero-shot prompts  \\
            ada & OpenAI & 350M & Base GPT-3 model from \citet{Brown20} \\
            davinci & OpenAI & 175B & Base GPT-3 model from \citet{Brown20} \\
            text-davinci-001 & OpenAI & 175B & Human-feedback model~\citep{ouyang2022training}; trained via supervised fine-tuning on human-written demonstrations. \\
            text-davinci-002 & OpenAI & 175B & Human-feedback model based on code-davinci-002~\citep{ouyang2022training}; trained via supervised fine-tuning on human-written demonstrations. \\
            text-davinci-003 & OpenAI & 175B & Improved version of text-davinci-002~\citep{ouyang2022training} \\
            \bottomrule
            \end{tabular}
    \end{small}
 \end{center}    
 \vskip -0.1in
 \end{table*}
 

Once we prompt a model with a given question, we simply evaluate the log probabilities that each of the answer choices is the next-token.
We then take these token log probabilities for each answer, exponentiate and then normalize them to get the model opinion distribution, i.e.,
$D_M = [e^{lp_A}, e^{lp_B}, ...,] / sum([e^{lp_A}, e^{lp_B}, ...,])$


Currently, OpenAI and AI21 limit the number of log probabilities they return via their API to 100 and 10 respectively.
Thus, if one of the option choices (say `A') is not in the set of returned log probabilities, we attempt to bound it as follows.
Let's say the model returns a set of $K$ (100 or 64) token-log probabilities pairs $\{t_k, lp_k\}$.
We compute the total assigned mass as $p_{assigned} = \sum_{k \in K} e^{lp_k}$. The remaining mass is thus $p_{missing} = 1 - M$. We also find $p_min = \min_{k in K}e^{lp_k}$, i.e., the minimum probability assigned to any of the $K$ token choices.
Then, we assigning the missing token `A' the probability $min(p_{missing}, p_{min})$. Note that this is an upper bound on the true probability mass the model assigns to token `A'.

As a baseline, we also consider a random model that chooses one of the answers choices per question at random.

\subsection{Metrics}
\label{app:metric}
To compute the Wasserstein distance between human and LM opinion distributions to a question, we must map the to options to a metric space.
To do so, we leverage the ordinal structure of the options (as provided by Pew surveys).
For instance, we would map the set of options `Strong Agree', `Agree', `Maybe', `Disagree' and `Strong Disagree' to the integers 1 through 5.
We follow this approach in most cases, with the exception being questions for which the penultimate option is non-ordinal. For instance, if the choices were `Very good', `Very bad', and `Neither good nor bad'. In this case, we map the answers to 1, 2 and 1.5 respectively.

\subsection{Temperature scaling}
\label{app:temp}

In Section~\ref{sec:rep}, we compare the model opinion distribution to a sharpened version of its human counterpart.
This sharpening makes the human opinion distribution collapse towards its dominant mode.
To do so, we use the standard temperature scaling approach from \citet{Guo17}.
We use a temperature of 1e-3 in our analysis, but find that our results are fairly robust to the choice of temperature.
